\documentclass[11pt,reqno]{amsart}
%==============================================================
%  Transport and tiling bounds for the affine plank problem on the triangle:
%  new sharp cases, a rigidity theorem, and a normalizer obstruction.
%
%  ACADEMIC-CORRECTNESS POLICY (binding): every Theorem/Proposition/Lemma is
%  either cited (with attribution) or proved in full here. Not-yet-established
%  statements live in Conjecture/Problem/Remark and are never used as proved.
%  SCOPE (honest): this studies the triangle as a CONCRETE body plus method.
%  It is NOT an attack on the general conjecture via Ambrus (whose reduction
%  targets simplices of dimension 2N-1; the triangle is never a target).
%==============================================================
\usepackage{amsmath,amssymb,amsthm,mathtools}
\usepackage[margin=1.1in]{geometry}
\usepackage[colorlinks=true,linkcolor=blue,citecolor=blue]{hyperref}
\theoremstyle{plain}
\newtheorem{theorem}{Theorem}[section]
\newtheorem{proposition}[theorem]{Proposition}
\newtheorem{lemma}[theorem]{Lemma}
\newtheorem{corollary}[theorem]{Corollary}
\theoremstyle{definition}
\newtheorem{conjecture}[theorem]{Conjecture}
\newtheorem{problem}[theorem]{Problem}
\theoremstyle{remark}
\newtheorem{remark}[theorem]{Remark}
\newcommand{\rw}{\operatorname{rw}}
\newcommand{\Leb}{\operatorname{Leb}}
\newcommand{\R}{\mathbb{R}}
\DeclareMathOperator{\wid}{w}
\DeclareMathOperator{\vol}{vol}

\title{Transport and tiling bounds for the affine plank problem on the triangle}
% AUTHOR BLOCK — decision pending (see 03-work-plan.md §5); fill before submission.
\author{[Author names]}
\address{[Affiliation]}
\email{[email]}
\date{Draft --- \today}

\begin{document}
\begin{abstract}
Bang's affine (relative-width) plank conjecture asserts that if a convex body $K$
is covered by planks $P_1,\dots,P_n$ then $\sum_i \rw_K(P_i)\ge 1$, where
$\rw_K(P)=\wid(P)/\wid_K(u_P)$ is the width of $P$ relative to the width of $K$ in
$P$'s direction. It is open for non-symmetric bodies, and open already for three
planks in the plane. We study the triangle as a concrete test body. We give (i) a
short single-measure proof of $\sum\rw\ge 1/d$ for $d\ge2$, whose per-plank estimate
is sharp at the simplex, so that no single-measure argument can improve on $1/d$;
(ii) sharp
$\sum\rw\ge 1$ for two directions and for facet-parallel families (all dimensions);
(iii) a Hunter-free proof of the ``three facets $+$ one arbitrary plank'' case; and,
as our main results, (iv) what is, to our knowledge, the first \emph{tight
non-facet} case with full rigidity:
three planks parallel to the medians of a triangle satisfy $\sum\rw\ge1$ with
equality if and only if each is the central third of its median; (v) a
\emph{weighted-perimeter theorem} extending the bound to a two-parameter family of
generic concurrent direction triples, via an explicit witness measure with edge
masses $1-2p_j$; and (vi) tightness and rigidity across that entire family: the
bound is attained at every member by an explicit covering, unique for its
directions --- a two-parameter family of sharp configurations. We isolate the
transport method behind (i)--(vi) as a
\emph{uniform-marginal witness measure} and give a dimension-agnostic necessary
condition for its existence (concurrence of the mid-level lines); on the
triangle the condition is also sufficient, giving a complete characterization
for three directions: a witness measure exists if and only if the mid-level
lines concur. Finally we
record a structural obstruction: the chord-lemma step in the sign argument of Bang,
which yields the chord-normalized bound $\sum \wid(P_i)/\ell_K(u_i)\ge1$, is sharp,
so that argument cannot be pushed past the longest chord: no normalizer strictly
between chord and width is reachable by it. We claim no proof of the conjecture.
\end{abstract}
\maketitle

\section{Introduction}\label{sec:intro}

A \emph{plank} of width $w$ is the region between two parallel hyperplanes at
distance $w$. For a convex body $K\subset\R^d$ and a plank $P$ with unit normal $u$,
the \emph{relative width} of $P$ is $\rw_K(P)=\wid(P)/\wid_K(u)$, where
$\wid_K(u)=\max_{x\in K}\langle u,x\rangle-\min_{x\in K}\langle u,x\rangle$. Bang's
affine plank conjecture~\cite{Bang51,Ambrus} states:
\begin{conjecture}[Affine plank problem; Bang 1951]\label{conj:bang}
If $K$ is covered by planks $P_1,\dots,P_n$, then $\sum_{i=1}^n \rw_K(P_i)\ge 1$.
\end{conjecture}
Ball~\cite{Ball91} proved Conjecture~\ref{conj:bang} for centrally symmetric $K$.
For general bodies it is open, and, as emphasised by Bakaev--Yehudayoff~\cite{BY26},
it is open \emph{already in the plane} --- only the two-plank case is settled there.
The best general lower bound is
$\sum\rw\ge 2/(1+\sqrt d)$~\cite{BY26}; on the equilateral triangle their method
gives $\sum\rw\ge 4\sqrt3-6\approx0.928$.

\smallskip
\noindent\textbf{Scope.}\ This paper studies the triangle as a concrete body and
develops the associated \emph{transport} (witness-measure) and \emph{tiling}
methods. It is \emph{not} an attack on Conjecture~\ref{conj:bang} through the Ambrus
reduction: that reduction sends a covering by $N$ planks to a simplex of dimension
$2N-1$, so the triangle (the $2$-simplex) is never a target
(Remark~\ref{rem:ambrus}). Our results are new sharp cases for a fixed body, a
rigidity theorem, a dimension-agnostic structural characterization, and an
obstruction; we prove nothing about the general conjecture.

\begin{remark}\label{rem:ambrus}
Ambrus's reduction~\cite{Ambrus} sends a covering by $N$ planks to a covering of a
simplex of dimension $2N-1$ (two points per plank). The target simplices thus have
odd dimension growing with $N$; the $2$-simplex (triangle) is never among them. So
proving the conjecture for the triangle does \emph{not} imply the planar case, and
our results are about the triangle as a concrete body, not a reduction of
Conjecture~\ref{conj:bang}.
\end{remark}

\subsection*{The triangle model}
By affine invariance take $\Delta=\{x\in\R^3:x_i\ge0,\ x_1+x_2+x_3=1\}$. A
normalized plank is $P=\{x\in\Delta: f(x)\in I\}$ with $f:\Delta\to[0,1]$ affine
onto $[0,1]$ and $I\subset[0,1]$ an interval; then $\rw_\Delta(P)=|I|$. Covering
means $\Delta\subset\bigcup_i P_i$. The restriction $I\subset[0,1]$ loses no
generality: replacing $I$ by $I\cap[0,1]$ leaves $P\cap\Delta$ unchanged and does
not increase $|I|$.

\subsection*{Results}
Theorem~\ref{thm:oned} (\S\ref{sec:measure}), Theorems~\ref{thm:twodir}
and~\ref{thm:facet} (\S\ref{sec:sharp}), Theorem~\ref{thm:3plus1}
(\S\ref{sec:3plus1}), and the centerpieces: Theorem~\ref{thm:median} with rigidity
(\S\ref{sec:median}); the weighted-perimeter Theorem~\ref{thm:wper}
(\S\ref{sec:concur}), which extends the sharp bound from the medians to a
two-parameter family of generic concurrent triples; and
Theorem~\ref{thm:tight}, which shows the bound is attained at every member of
that family by an explicit unique tight covering.
Proposition~\ref{prop:concur} (\S\ref{sec:concur}) is the dimension-agnostic
necessary condition, Corollary~\ref{cor:iff} the resulting characterization for
cyclic triples, and Theorem~\ref{thm:char} the full three-direction
characterization on the triangle; Proposition~\ref{thm:nogo} (\S\ref{sec:nogo})
is the normalizer obstruction. Appendix~\ref{app:tilings} classifies the
boundary tilings of cyclic triples (Proposition~\ref{prop:tilings}), the
combinatorial engine behind both rigidity results.

\section{The single-measure bound $\sum\rw\ge 1/d$}\label{sec:measure}
\begin{theorem}\label{thm:oned}
Let $d\ge2$. If $K\subset\R^d$ is covered by planks $P_1,\dots,P_n$, then
$\sum_i\rw_K(P_i)\ge 1/d$. The per-plank estimate $\mu(P)\le d\,\rw_K(P)$ used in the
proof is sharp at the simplex; hence no single-measure argument can improve the
constant $1/d$. (For $d=1$ the conjecture itself is trivial: intervals covering an
interval satisfy $\sum\rw\ge1$.)
\end{theorem}
\begin{proof}
Let $\mu$ be the uniform probability measure on $K$. Fix $u$, normalize
$\langle u,\cdot\rangle\in[0,w]$, $w=\wid_K(u)$. The marginal density is
$f(t)=V(t)/\vol(K)$ with $V(t)=\vol_{d-1}(K\cap\{\langle u,\cdot\rangle=t\})$; by
Brunn--Minkowski $g:=V^{1/(d-1)}$ is concave on $[0,w]$. If $g$ peaks at $t^\ast$
with value $g^\ast$, then $g\ge g^\ast h$ for the tent $h$ ($h(0)=h(w)=0$,
$h(t^\ast)=1$), so $f=g^{d-1}\ge(\max f)\,h^{d-1}$ and
$1=\int f\ge(\max f)\int_0^w h^{d-1}=(\max f)\,w/d$, giving $\max f\le d/w$. Hence
$\mu(P)=\int_I f\le d\,\rw_K(P)$ for each plank, and
$1=\mu(K)\le\sum_i\mu(P_i)\le d\sum_i\rw_K(P_i)$. For sharpness of the per-plank
estimate: equality in $\max f\le d/w$ forces $V\propto h^{d-1}$ (conical sections),
realized by the simplex $K=\Delta^d$ with $u$ a vertex--facet direction, where a
thin plank at the facet has $\mu(P)=(d-o(1))\rw_K(P)$. Thus for the simplex no
single probability measure can certify a constant larger than $1/d$ by the union
bound; this does \emph{not} assert any covering of the simplex with
$\sum\rw$ near $1/d$ (indeed Theorem~\ref{thm:facet} gives $\sum\rw\ge1$ for
facet-parallel coverings).
\end{proof}
\begin{remark}
The bound cannot exceed $1/d$ by a single measure: uniform marginals in all
directions would force $\mathbb E[x_k]=1/2$ against $\sum x_k=1$. This is Gardner's
obstruction~\cite{Gardner88}: for the triangle there is no relative-width measure
for the three side-directions.
\end{remark}

\section{Sharp cases: two directions and facet-parallel families}\label{sec:sharp}
\begin{theorem}\label{thm:twodir}
If every plank covering the triangle is parallel to one of two fixed
\emph{non-parallel} directions, then $\sum\rw\ge1$, and this is sharp. (If the two
directions are parallel the statement is trivial: all planks share one normalized
coordinate $f$, the intervals $f(P_i)$ cover $[0,1]$, and $\sum\rw=\sum|I_i|\ge1$.)
\end{theorem}
\begin{proof}
Let $f_u,f_v:\Delta\to[0,1]$ be the two affine coordinates; since the directions are
non-parallel, $\Phi=(f_u,f_v)$ is an
affine bijection onto a triangle $T\subset[0,1]^2$ with full projections. By
\cite[Theorem~1]{Gardner88}, which constructs a relative-width measure for any
\emph{two} directions on any planar convex body,
there is a probability measure $\nu$ on $\Delta$ with $f_{u\#}\nu=f_{v\#}\nu=\Leb$.
Then $\nu(P)=|I|=\rw(P)$ for each plank, and
$1=\nu(\Delta)\le\sum\nu(P_i)=\sum\rw(P_i)$.
\end{proof}
\begin{theorem}\label{thm:facet}
Planks parallel to the facets of a simplex $\Delta^d$, in any number, that cover
$\Delta^d$ satisfy $\sum\rw\ge1$, and this is sharp.
\end{theorem}
\begin{proof}
Write $\Delta^d=\{\lambda\in[0,1]^{d+1}:\sum_k\lambda_k=1\}$ in barycentric
coordinates; a plank parallel to facet $k$ is a coordinate slab
$\{\lambda_k\in I\}$. Grouping the planks by facet, let
$U_k=\bigcup_{c(i)=k}I_i\subset[0,1]$ be the forbidden set on facet $k$, so
$\sum_k|U_k|\le\sum_i|I_i|=:R$, and set $C_k=[0,1]\setminus U_k$. A point
$\lambda\in\Delta^d$ is uncovered iff $\lambda_k\in C_k$ for every $k$; since
$C_k\subset[0,1]$, an uncovered point exists \emph{iff}
$1\in C_1+\dots+C_{d+1}$ (Minkowski sum).

Assume $R<1$; we produce an uncovered point. Each $C_k$ is nonempty compact
($|C_k|=1-|U_k|\ge1-R>0$). Put $D=C_1+\dots+C_d$. By the one-dimensional
Brunn--Minkowski inequality $|A+B|\ge|A|+|B|$ (for nonempty compact $A,B\subset\R$,
iterated), $|D|\ge\sum_{k=1}^d|C_k|=d-\sum_{k=1}^d|U_k|$; as $D\subseteq[0,d]$,
$|[0,1]\setminus D|\le|[0,d]\setminus D|\le\sum_{k=1}^d|U_k|$, whence
$|D\cap[0,1]|\ge1-\sum_{k=1}^d|U_k|$. Also $(1-C_{d+1})\cap[0,1]$ has measure
$|C_{d+1}|=1-|U_{d+1}|$. Both sets lie in $[0,1]$ and their measures sum to
$2-\sum_k|U_k|\ge2-R>1$, so they intersect: there is $t\in D$ with $1-t\in C_{d+1}$.
Then $t+(1-t)=1\in C_1+\dots+C_{d+1}$, i.e.\ an uncovered $\lambda$ exists.
Contrapositively, covering forces $R\ge1$. Sharpness: the facet tiling
$I_k=[0,\tfrac1{d+1}]$ gives $\sum\rw=1$.
\end{proof}
\begin{remark}
The edge-by-edge argument does \emph{not} prove this: on an edge only two
barycentric coordinates vary, so no single edge can force $\sum_k|U_k|\ge1$ across
all $d+1$ facets. The Minkowski-sum reduction above is the correct route; note that
no witness \emph{measure} can prove it either, since facet directions violate the
concurrence condition of Proposition~\ref{prop:concur}.
\end{remark}

\section{Three facets plus one arbitrary plank}\label{sec:3plus1}
\begin{theorem}\label{thm:3plus1}
Let $P_i=\{\lambda_i\in I_i\}$ ($i=1,2,3$) be facet-parallel with $a_i=|I_i|$, and
let $Q=\{f\in J\}$, $t=|J|$, with $f:\Delta\to[0,1]$ affine normalized. If
$\Delta\subset P_1\cup P_2\cup P_3\cup Q$ then $a_1+a_2+a_3+t\ge1$.
\end{theorem}
\begin{proof}
Choose an active edge of $f$; relabelling, $f(v_j)=0$, $f(v_k)=1$, and let $v_i$ be
the opposite vertex, $\theta=f(v_i)$. For $s\in[0,1]$ consider the fibre
$F_s=\{\lambda_i=s\}$, parametrized by $u=\lambda_k\in[0,1-s]$, so
$\lambda_j=1-s-u$ and $f|_{F_s}(u)=\theta s+u$. On $F_s$: $P_i\cap F_s=F_s$ if
$s\in I_i$ (else $\varnothing$); $P_j\cap F_s$, $P_k\cap F_s$, $Q\cap F_s$ have
lengths $\le a_j,a_k,t$. Put $B=a_j+a_k+t$. If $s\notin I_i$ and $F_s$ is covered,
its length $1-s\le B$. If $B\ge1$ we are done; else $[0,1-B)\subset I_i$, so
$a_i\ge1-B$. In both cases $a_1+a_2+a_3+t=a_i+B\ge1$.
\end{proof}

\section{The median theorem and its rigidity}\label{sec:median}
Let $m_1,m_2,m_3$ be the medians of $\Delta$ as normalized affine forms: with
$A,B,C$ the vertices,
\[ m_1=(\tfrac12,0,1),\quad m_2=(0,1,\tfrac12),\quad m_3=(1,\tfrac12,0), \]
(vertex values). One has $m_1+m_2+m_3\equiv\tfrac32$. A \emph{median plank} is
$P_i=\{m_i\in I_i\}$, $I_i=[l_i,h_i]$, $r_i=|I_i|>0$.
\begin{theorem}\label{thm:median}
If three median planks cover $\Delta$, then $r_1+r_2+r_3\ge1$, with equality if and
only if $I_1=I_2=I_3=[\tfrac13,\tfrac23]$.
\end{theorem}
To our knowledge this is the first tight \emph{non-facet} case of the three-plank
problem with full rigidity: the sharp cases recorded in the survey
\cite{Verreault} (and in \cite{Gardner88,BY26}) concern facet/coordinate directions,
two directions, or symmetric bodies, and we are not aware of a comparable tilted
tight instance in the literature.
\begin{proof}[Proof of the bound]
Let $\mu$ be the perimeter measure: uniform in the parameter on each of the three
edges, mass $\tfrac13$ each. A direct computation shows each pushforward
$m_{i\#}\mu=\Leb[0,1]$. Indeed for $m_1$: on $AB$, $m_1$ ranges over $[0,\tfrac12]$
with Jacobian $2$ (density $\tfrac13\cdot2$); on $CA$, over $[\tfrac12,1]$ likewise;
on $BC$, over $[0,1]$ with Jacobian $1$ (density $\tfrac13$); the total density is
$1$ on $[0,1]$. Hence $\mu(P_i)=\Leb(I_i)=r_i$,
and $1=\mu(\Delta)\le\sum\mu(P_i)\le\sum r_i$.
\end{proof}
\begin{proof}[Proof of rigidity]
Equality in $1\le\sum\mu(P_i)$ forces the $P_i$ to
partition $\mu$ a.e.; since $\mu$ lives on the edges, \emph{on each edge the three
traces tile $[0,1]$}. We claim there are exactly three such edge-tilings, and only
one covers the interior.

\emph{Three edge-tilings.} On each edge each median is affine, so the three traces
are explicit intervals depending piecewise-linearly on $(l_i,h_i)$ through
$a_i=|I_i\cap[0,\tfrac12]|$ and $b_i=|I_i\cap[\tfrac12,1]|$. Classifying each $I_i$ by
crossing type (contained in $[0,\tfrac12]$, crossing $\tfrac12$, or contained in
$[\tfrac12,1]$) and each edge by the left-to-right order of its traces, the abutting
(tiling) conditions become a finite family of linear systems in $(l_i,h_i)$.

\begin{lemma}\label{lem:three}
The edge-tiling conditions have exactly three solutions with $0\le l_i<h_i\le1$:
$[\tfrac13,\tfrac23]^3$, $[0,\tfrac13]^3$, and $[\tfrac23,1]^3$.
\end{lemma}
This is the case $\tau_i=\tfrac12$ of Proposition~\ref{prop:tilings}, proved in
full in Appendix~\ref{app:tilings} for arbitrary cyclic triples: the symmetries
of the configuration reduce the a priori $27$ crossing patterns to seven
classes, five of which are impossible and two of which lead to cyclic linear
systems with unique solutions. The lemma has also been verified by three
independent exact computer enumerations, archived as ancillary files.

\emph{Interior witness.} The centroid $G$ has $m(G)=(\tfrac12,\tfrac12,\tfrac12)$.
Since $\tfrac12\notin[0,\tfrac13]\cup[\tfrac23,1]$, the tilings $[0,\tfrac13]^3$ and
$[\tfrac23,1]^3$ leave $G$ uncovered; they tile the boundary but are not coverings
of $\Delta$. The tiling $[\tfrac13,\tfrac23]^3$ covers $\Delta$ and has
$\sum r_i=1$. Hence the unique covering with $\sum r_i=1$ is $[\tfrac13,\tfrac23]^3$.
\end{proof}

\section{A dimension-agnostic necessary condition for transport}\label{sec:concur}
The proofs of Theorems~\ref{thm:twodir} and~\ref{thm:median} share one device: a
probability measure whose pushforward under each direction is uniform. Which
direction-tuples admit one?
\begin{proposition}\label{prop:concur}
Let $u_1,\dots,u_{d+1}$ be normalized affine forms on $\Delta^d$ with vertex-value
matrix $V$ invertible. If a probability measure $\mu$ on $\Delta^d$ has
$u_{i\#}\mu=\Leb[0,1]$ for all $i$, then its barycenter $p\in\Delta^d$ satisfies
$Vp=\tfrac12\mathbf 1$; equivalently the mid-level hyperplanes $\{u_i=\tfrac12\}$
concur at $p$, which requires $\mathbf 1^\top V^{-1}\mathbf 1=2$ and
$V^{-1}\mathbf 1\ge0$.
\end{proposition}
\begin{proof}
$\tfrac12=\mathbb E_\mu[u_i]=\sum_j V_{ij}\mathbb E_\mu[x_j]=(Vp)_i$, and
$p\in\Delta^d$ gives $\mathbf 1^\top p=1$, $p\ge0$; substitute $p=\tfrac12V^{-1}\mathbf1$.
\end{proof}
The three medians satisfy this with $p$ the centroid; the three facets do not
($\mathbf 1^\top V^{-1}\mathbf 1=3\ne2$), recovering Gardner's obstruction. The
concurrence family is two-dimensional (parametrized by $p$): a normalized direction
has vertex values $\{0,\tau,1\}$, and $Vp=\tfrac12\mathbf1$ ties each $\tau_i$ to $p$.
The medians are the case $p=$ centroid, forcing $\tau_i=\tfrac12$. Among these, the
explicit perimeter measure singles out the medians:
\begin{proposition}\label{prop:tauhalf}
For a direction $u$ with vertex values $\{0,\tau,1\}$, the perimeter measure $\mu$
(uniform on each edge, mass $\tfrac13$) satisfies $u_\#\mu=\Leb[0,1]$ if and only if
$\tau=\tfrac12$.
\end{proposition}
\begin{proof}
On the three edges $u$ is affine between the pairs $\{0,\tau\},\{\tau,1\},\{0,1\}$,
so $u_\#\mu$ has density $\tfrac1{3\tau}+\tfrac13$ on $[0,\tau]$ and
$\tfrac1{3(1-\tau)}+\tfrac13$ on $[\tau,1]$. Both equal $1$ iff $\tau=\tfrac12$.
\end{proof}
Uniform edge-weights thus single out the medians. Allowing \emph{unequal} edge
masses, however, the construction covers a full two-parameter family:
\begin{theorem}[weighted perimeter]\label{thm:wper}
Let $p=(p_A,p_B,p_C)$ be interior with $\max_j p_j<\tfrac12$, and let
$u_1=(\tau_1,0,1)$, $u_2=(0,1,\tau_2)$, $u_3=(1,\tau_3,0)$ (vertex values at
$A,B,C$) be the cyclic directions concurrent at $p$, i.e.\
$\tau_1=\frac{1/2-p_C}{p_A}$, $\tau_2=\frac{1/2-p_B}{p_C}$,
$\tau_3=\frac{1/2-p_A}{p_B}$. Let $\mu_p$ be uniform on each edge with masses
$w_{AB}=1-2p_C$, $w_{BC}=1-2p_A$, $w_{CA}=1-2p_B$ (summing to $1$). Then
$u_{i\#}\mu_p=\Leb[0,1]$ for $i=1,2,3$; consequently every covering of $\Delta$ by
planks in these three directions satisfies $\sum\rw\ge1$.
\end{theorem}
\begin{proof}
Each $u_i$ is affine of full range $[0,1]$ on the edge joining its $0$- and
$1$-vertices ($u_1$ on $BC$, $u_2$ on $AB$, $u_3$ on $CA$) and affine of range
$[0,\tau_i]$, resp.\ $[\tau_i,1]$, on the other two. The pushforward of a uniform
edge mass $w_e$ through an affine map of range length $L$ is uniform with density
$w_e/L$. Hence $u_{1\#}\mu_p$ has density $w_{BC}+w_{AB}/\tau_1$ on $[0,\tau_1]$ and
$w_{BC}+w_{CA}/(1-\tau_1)$ on $[\tau_1,1]$, and similarly for $u_2,u_3$: six
equations. Using $p_A+p_B+p_C=1$, each composite term collapses by the identity
$(1-2x)\,p/(\tfrac12-x)=2p$: e.g.\
$w_{AB}/\tau_1=(1-2p_C)p_A/(\tfrac12-p_C)=2p_A$, so
$w_{BC}+w_{AB}/\tau_1=(1-2p_A)+2p_A=1$; likewise
$1-\tau_1=\frac{1/2-p_B}{p_A}$ gives $w_{CA}/(1-\tau_1)=2p_A$. All six densities
equal $1$, so each marginal is $\Leb[0,1]$, and the transport argument of
Theorem~\ref{thm:median} applies verbatim. Finally $\tau_i\in(0,1)$ for all $i$ iff
$\max_jp_j<\tfrac12$ iff all three masses are positive, so every valid cyclic triple
admits $\mu_p$.
\end{proof}
Taking $p$ the centroid recovers Theorem~\ref{thm:median}'s bound ($\tau_i=\tfrac12$,
equal masses). For $p$ off the centroid the directions are generic non-median
triples; e.g.\ $p=(\tfrac9{20},\tfrac3{10},\tfrac14)$ gives
$\tau=(\tfrac59,\tfrac45,\tfrac16)$ with masses $(\tfrac12,\tfrac1{10},\tfrac25)$.
Theorem~\ref{thm:wper} thus proves the bound of Conjecture~\ref{conj:bang} for a
two-parameter family of generic direction triples on the triangle, answering the
finite-direction existence question of~\cite{Gardner88} affirmatively for the cyclic
concurrent family.

Combining necessity (Proposition~\ref{prop:concur}) with the explicit construction
(Theorem~\ref{thm:wper}) yields the first \emph{characterization} in this setting:
\begin{corollary}\label{cor:iff}
Call a direction triple \emph{cyclic} if its vertex values are
$u_1=(\tau_1,0,1)$, $u_2=(0,1,\tau_2)$, $u_3=(1,\tau_3,0)$ with
$\tau_i\in(0,1)$. A cyclic triple admits a uniform-marginal witness measure if and
only if the three mid-level lines $\{u_i=\tfrac12\}$ concur. In that case the
concurrence point $p$ automatically lies in the open medial triangle
($p_j>0$ and $p_j<\tfrac12$ for all $j$), and the weighted-perimeter measure
$\mu_p$ of Theorem~\ref{thm:wper} is a witness.
\end{corollary}
\begin{proof}
The vertex-value matrix has $\det V=-(1+\tau_1\tau_2\tau_3)\ne0$, so
Proposition~\ref{prop:concur} applies and gives necessity. Conversely, if the lines
concur at $p$ (a point of the affine plane of $\Delta$, so $\sum_jp_j=1$), then $p$
is the unique solution of $Vp=\tfrac12\mathbf1$, which in closed form reads
$2(1+\tau_1\tau_2\tau_3)\,p_A=1-\tau_3(1-\tau_2)$ and cyclically; each right-hand
side is positive for $\tau_i\in(0,1)$, so $p_j>0$, and the identities
$\tfrac12-p_C=\tau_1p_A$, $\tfrac12-p_A=\tau_3p_B$, $\tfrac12-p_B=\tau_2p_C$ give
$p_j<\tfrac12$. Theorem~\ref{thm:wper} then supplies $\mu_p$.
\end{proof}
\begin{remark}
The \emph{anti-cyclic} pattern (all $0$- and $1$-roles exchanged) reduces to the
cyclic one by a reflection of $\Delta$ (a transposition of two vertices), which
maps the medial triangle to itself; Theorem~\ref{thm:wper},
Corollary~\ref{cor:iff} and Theorem~\ref{thm:tight} below hold for it verbatim.
All remaining role assignments are classified by Theorem~\ref{thm:char} below.
\end{remark}

The bound of Theorem~\ref{thm:wper} is in fact \emph{attained} at every member
of the family, and the tight covering is unique --- a two-parameter family of
sharp configurations extending the median rigidity of \S\ref{sec:median}.

\begin{theorem}[tightness and rigidity of the cyclic family]\label{thm:tight}
Let $p$ and $u_1,u_2,u_3$ be as in Theorem~\ref{thm:wper}. Write
$\alpha=1-2p_A$, $\beta=1-2p_B$, $\gamma=1-2p_C$ (the edge masses of $\mu_p$;
$\alpha,\beta,\gamma>0$, $\alpha+\beta+\gamma=1$) and
$S=\alpha\beta+\beta\gamma+\gamma\alpha$. Then the planks $P_i=\{u_i\in I_i\}$
with
\[
I_1=\Bigl[\tfrac{\alpha\gamma}S,\,1-\tfrac{\alpha\beta}S\Bigr],\qquad
I_2=\Bigl[\tfrac{\beta\gamma}S,\,1-\tfrac{\alpha\gamma}S\Bigr],\qquad
I_3=\Bigl[\tfrac{\alpha\beta}S,\,1-\tfrac{\beta\gamma}S\Bigr]
\]
cover $\Delta$, with $\sum\rw
=\tfrac{\beta\gamma}S+\tfrac{\alpha\beta}S+\tfrac{\alpha\gamma}S=1$. Moreover
this is the unique covering of $\Delta$ with $\sum\rw=1$ by three planks of
positive width, one in each of the directions $u_1,u_2,u_3$. At $p$ the centroid
($\alpha=\beta=\gamma=\tfrac13$) it is $[\tfrac13,\tfrac23]^3$, recovering
Theorem~\ref{thm:median}.
\end{theorem}

\begin{proof}
Write $l=(l_1,l_2,l_3)=(\alpha\gamma,\beta\gamma,\alpha\beta)/S$ and
$h=(h_1,h_2,h_3)=\mathbf1-(\alpha\beta,\alpha\gamma,\beta\gamma)/S$, and note
$\tau_1=\frac{\gamma}{1-\alpha}$, $\tau_2=\frac{\beta}{1-\gamma}$,
$\tau_3=\frac{\alpha}{1-\beta}$ from the formulas of Theorem~\ref{thm:wper}. A
direct substitution shows $(l,h)=(\lambda,\rho)$, the first solution of
Proposition~\ref{prop:tilings}; in particular $l_i<\tau_i<h_i$, and the traces
of $P_1,P_2,P_3$ tile all three edges, so $\partial\Delta\subset\bigcup P_i$.

\emph{Step 1: an affine identity.} Let $a_1=\alpha(1-\alpha)$,
$a_2=\gamma(1-\gamma)$, $a_3=\beta(1-\beta)$ and
$\Lambda=a_1u_1+a_2u_2+a_3u_3$. At the vertices:
$\Lambda(A)=a_1\tau_1+a_3=\alpha\gamma+\beta(\alpha+\gamma)=S$,
$\Lambda(B)=a_2+a_3\tau_3=\gamma(\alpha+\beta)+\alpha\beta=S$,
$\Lambda(C)=a_1+a_2\tau_2=\alpha(\beta+\gamma)+\beta\gamma=S$. An affine
function agreeing at three affinely independent points is constant:
$\Lambda\equiv S$ on the plane of $\Delta$. Moreover
$\sum_ia_i=1-(\alpha^2+\beta^2+\gamma^2)=2S$, and expanding (using
$1-\alpha=\beta+\gamma$ etc.)
\[
\textstyle\sum_i a_il_i
=\frac{\alpha^2\gamma(\beta+\gamma)+\beta\gamma^2(\alpha+\beta)
+\alpha\beta^2(\gamma+\alpha)}{S}
=\frac{(\alpha^2\beta^2+\beta^2\gamma^2+\gamma^2\alpha^2)+\alpha\beta\gamma}{S}
=\frac{S^2-\alpha\beta\gamma}{S},
\]
since $S^2=\sum\alpha^2\beta^2+2\alpha\beta\gamma$; the same expansion gives
$\sum_ia_i(1-h_i)=\frac{S^2-\alpha\beta\gamma}S$, whence
$\sum_ia_ih_i=2S-\frac{S^2-\alpha\beta\gamma}S=S+\frac{\alpha\beta\gamma}S$.
Thus
\begin{equation}\label{eq:margin}
\textstyle\sum_i a_il_i=S-\frac{\alpha\beta\gamma}S<S<
S+\frac{\alpha\beta\gamma}S=\sum_i a_ih_i .
\end{equation}

\emph{Step 2: the planks cover $\Delta$.} Suppose not; the uncovered set $U$ is
relatively open (planks are closed) and misses $\partial\Delta$, and it
decomposes as $U=\bigcup_\sigma P_\sigma$ over the eight sign patterns
$\sigma\in\{\mathrm{lo},\mathrm{hi}\}^3$, where $P_\sigma=H_\sigma\cap
\operatorname{int}\Delta$ and $H_\sigma$ is the open convex set
$\{u_i<l_i\ (\sigma_i=\mathrm{lo}),\ u_i>h_i\ (\sigma_i=\mathrm{hi})\}$; write
$c_i$ for the active bound ($l_i$ or $h_i$). We show each $P_\sigma$ is empty.

First, $\overline{P_\sigma}\cap\partial\Delta$ is finite: such a point $x$
satisfies the closed pattern constraints, and is covered by some $P_j$ (the
boundary is tiled), forcing $u_j(x)=c_j$; each $u_j$ is nonconstant on each
edge (its vertex values there are distinct), so $x$ lies in a set of at most
nine points. Next, $H_\sigma\subseteq\overline\Delta$ whenever
$P_\sigma\ne\varnothing$: otherwise pick $y\in H_\sigma\setminus\overline\Delta$
and a disk $B\subset P_\sigma$; for each $z\in B$ the segment $[z,y]\subset
H_\sigma$ crosses $\partial\Delta$ at a point of
$\overline{P_\sigma}\cap\partial\Delta$, and segments through the fixed $y$
meeting a common crossing point have collinear $z$, so a disk of $z$'s produces
infinitely many crossing points --- contradicting finiteness. In particular
$H_\sigma$ is bounded.

No two of the lines $\ell_i=\{u_i=c_i\}$ are parallel: $u_i\parallel u_j$ would
force $\tau_1(1-\tau_2)=1$, $\tau_2(1-\tau_3)=1$ or $\tau_3(1-\tau_1)=1$,
impossible for $\tau\in(0,1)^3$. If the three lines are concurrent, $H_\sigma$
is an intersection of three open half-planes whose boundaries pass through one
point, hence empty or an unbounded open cone; boundedness gives
$H_\sigma=\varnothing$. Otherwise $H_\sigma$, if nonempty, is the open triangle
with vertices $v_{ij}=\ell_i\cap\ell_j\in\overline{H_\sigma}
\subseteq\overline\Delta$, and each $v_{ij}$ satisfies the closed constraint in
the third index $k$. Since $\Lambda\equiv S$,
$u_k(v_{ij})=(S-a_ic_i-a_jc_j)/a_k$, so that constraint reads $S\le\sum_ma_mc_m$
if $\sigma_k=\mathrm{lo}$, and $S\ge\sum_ma_mc_m$ if $\sigma_k=\mathrm{hi}$. For
a \emph{mixed} pattern both directions occur, forcing $S=\sum_ma_mc_m$; but then
$u_k(v_{ij})=c_k$, i.e.\ the three lines concur --- excluded. For
$\sigma=(\mathrm{lo},\mathrm{lo},\mathrm{lo})$ we get $S\le\sum a_il_i$, and for
$\sigma=(\mathrm{hi},\mathrm{hi},\mathrm{hi})$, $S\ge\sum a_ih_i$ --- both
contradict \eqref{eq:margin}. Hence $U=\varnothing$: the planks cover, and
$\sum\rw=1$.

\emph{Step 3: uniqueness.} Let $(I_i')$ be any covering with all
$r_i'=|I_i'|>0$ and $\sum r_i'=1$. Since $u_{i\#}\mu_p=\Leb[0,1]$
(Theorem~\ref{thm:wper}), $\mu_p(P_i')=r_i'$, and
$1=\mu_p(\Delta)\le\sum\mu_p(P_i')=1$ forces the planks to partition $\mu_p$
a.e.; as $\mu_p$ has positive uniform density on each edge, the traces tile all
three edges. By Proposition~\ref{prop:tilings}, $(I_i')$ is $(\,[\lambda_i,
\rho_i]\,)=(I_i)$, $(\,[0,\lambda_i]\,)$ or $(\,[\rho_i,1]\,)$. The last two do
not cover $\Delta$: set $\delta=\frac{\alpha\beta\gamma}{2S^2}$ and let
$x^\ast$ be the unique point of the plane with $u_i(x^\ast)=\lambda_i+\delta$
for all $i$ (consistent with $\Lambda\equiv S$ by \eqref{eq:margin} and
$\sum a_i=2S$; unique since $V$ is invertible). Its barycentric coordinates are
\[
x^\ast=\tfrac1{2S^2}\Bigl(\alpha\beta(1-\alpha)(1+\alpha-\beta),\;
\beta\gamma(1-\beta)(\alpha+2\beta),\;
\alpha\gamma(\alpha+\beta)(2-2\alpha-\beta)\Bigr),
\]
all positive, so $x^\ast\in\operatorname{int}\Delta$ and
$u_i(x^\ast)>\lambda_i$ for every $i$: $x^\ast$ is uncovered by
$(\,[0,\lambda_i]\,)$. Symmetrically the point $y^\ast$ with
$u_i(y^\ast)=\rho_i-\delta$, with coordinates
$\frac1{2S^2}\bigl(\alpha\gamma(1-\alpha)(2\alpha+\beta),\,
\alpha\beta(1-\beta)(1+\beta-\alpha),\,
\beta\gamma(\alpha+\beta)(2-\alpha-2\beta)\bigr)$, is uncovered by
$(\,[\rho_i,1]\,)$. Hence $(I_i')=(I_i)$.
\end{proof}

\begin{remark}
At the centroid, $\delta=\tfrac16$ and both witnesses degenerate to the point
with $u_i\equiv\tfrac12$ --- the centroid itself, recovering the interior
witness used in the proof of Theorem~\ref{thm:median}. The identities of Step~1
and the positivity of the witness coordinates, as well as exact coverage checks
at several family points, are verified symbolically in the archived ancillary
script \texttt{cyclic\_family\_tight\_rigidity.py}.
\end{remark}

Finally, the restriction to cyclic role assignments can be removed altogether.
The plank problem is invariant under the \emph{flip} of a normalized direction,
$u\mapsto1-u$, $I\mapsto1-I$ (the planks, their widths, and uniformity of
marginals are unchanged, and each mid-level line is fixed). Call the edge
joining the $0$- and $1$-vertices of a direction its \emph{full edge}. Flips
reduce every role assignment to one of three shapes, and the existence question
is settled in each:

\begin{theorem}[three-direction characterization]\label{thm:char}
Let $u_1,u_2,u_3$ be normalized directions on $\Delta$, pairwise non-parallel.
A probability measure $\mu$ on $\Delta$ with $u_{i\#}\mu=\Leb[0,1]$ for
$i=1,2,3$ exists if and only if the three mid-level lines $\{u_i=\tfrac12\}$
have a common point $p$. In that case $\sum\rw\ge1$ for every covering of
$\Delta$ by planks parallel to these directions, and exactly one of the
following holds:
\begin{itemize}
\item[(a)] the full edges are pairwise distinct: up to flips the triple is
cyclic, $p$ lies in the open medial triangle, and Theorems~\ref{thm:wper}
and~\ref{thm:tight} apply verbatim (explicit witness $\mu_p$; unique tight
covering);
\item[(b)] all three directions have the same full edge $e$: then $p$ is the
midpoint of $e$, the uniform measure on $e$ is a witness, and the bound follows
from covering $e$ alone (the trivial one-dimensional case recorded in
Theorem~\ref{thm:oned}).
\end{itemize}
\end{theorem}

\begin{proof}
\emph{Necessity.} If $\mu$ exists, its barycenter $p\in\Delta$ satisfies
$u_i(p)=\mathbb E_\mu[u_i]=\tfrac12$ for every $i$ (no invertibility of $V$ is
needed), so the mid-lines concur at $p$.

\emph{Classification.} (1) If the three full edges are distinct, flipping each
direction so that its $(0,1)$-pair is oriented cyclically makes the triple
cyclic; by Corollary~\ref{cor:iff}, concurrence then places $p$ in the open
medial triangle and yields $\mu_p$, and Theorem~\ref{thm:tight} the unique
tight covering. (2) Suppose exactly two directions, say $u_i$ and $u_j$, share
their full edge $e$, and let $x_v$ be the barycentric coordinate of the vertex
$v$ opposite $e$. After a flip, $u_i$ and $u_j$ agree on $e$, so
$u_i-u_j=(\tau_i-\tau_j)x_v$ with $\tau_i\ne\tau_j$ (equality would make them
parallel). Hence $\{u_i=\tfrac12\}\cap\{u_j=\tfrac12\}\subset\{x_v=0\}$, the
line through $e$, on which both mid-lines pass only through the midpoint of
$e$; but the third direction is not full on $e$, its values at the endpoints of
$e$ being $\{0,\tau_k\}$ or $\{\tau_k,1\}$, so its value at the midpoint is
$\tfrac{\tau_k}2$ or $\tfrac{1+\tau_k}2$, never $\tfrac12$. Thus the mid-lines
of such a triple \emph{never} concur; consistently, no witness measure exists
($\mathbb E_\mu[u_i]=\mathbb E_\mu[u_j]=\tfrac12$ would force
$\mathbb E_\mu[x_v]=0$, hence $\operatorname{supp}\mu\subset e$, on which $u_k$
is not onto $[0,1]$). This shape contributes to neither side of the
equivalence. (3) If all three directions share the full edge $e$, then after
flips each restricts to $e$ as the identity parameter, so all mid-lines pass
through the midpoint of $e$ (concurrence always holds), and the uniform
probability measure on $e$ has all three pushforwards equal to $\Leb[0,1]$.
For the bound, a covering of $\Delta$ covers $e$, on which the $i$-th plank
traces an interval of length $r_i$; hence $\sum r_i\ge1$.
\end{proof}

\begin{remark}
Theorem~\ref{thm:char} answers the existence question of~\cite{Gardner88}
completely for three directions on the triangle. In particular the open case of
Conjecture~\ref{conj:bang} on the triangle is reduced to \emph{non-concurrent}
triples, for which no single witness measure can exist.
\end{remark}

\smallskip
The concurrence condition is dimension-agnostic: the cyclic medians of $\Delta^d$
($m_i(e_i)=\tfrac12$, $m_i(e_{i+1})=0$, $m_i(e_{i-1})=1$, else $\tfrac12$) have
row-sums $\tfrac{d+1}2$, hence concur at the centroid in every dimension. Whether a
witness measure then exists for $d\ge3$ is open (the perimeter measure is special to
$d=2$); this is a higher-dimensional instance of the finite-direction question left
open by Gardner~\cite{Gardner88}.

\section{A normalizer obstruction}\label{sec:nogo}
For a direction $u$ let $\ell_K(u)$ be the length of the longest chord of $K$
parallel to $u$; always $\ell_K(u)\le\wid_K(u)$. Bakaev--Yehudayoff's chord bound
$\sum_i\wid(P_i)/\ell_K(u_i)\ge1$~\cite{BY26} follows from Bang's sign argument
via the chord lemma. Consider interpolating the normalizer,
$N_c=(1-c)\ell+c\,\wid$, $c\in[0,1]$: $N_c=\ell$ at $c=0$ (chord bound) and
$N_c=\wid$ at $c=1$ (Conjecture~\ref{conj:bang}).
\begin{proposition}\label{thm:nogo}
The chord-lemma step is sharp: for $u$ with $\|u\|=a/2$, the body
$(K-u)\cap(K+u)$ contains a homothet of $K$ of ratio $(\ell-a)/\ell$
\cite[Lem.~2.3]{Verreault}, and this ratio is attained (by a segment). Consequently
the iterated chord-lemma placement of a Bang-system witness $x_0+\sum_j\pm u_j$
proves $\sum_i\wid(P_i)/N(u_i)\ge1$ for $N=\ell$ and for no direction-normalizer
$N>\ell$; in particular it does not establish $\sum_i\wid(P_i)/N_c(u_i)\ge1$ for
any $c>0$.
\end{proposition}
\begin{proof}
The placement argument runs as follows: iterating the chord lemma over the
directions $u_j$, the common intersection $\bigcap_j\bigl((K-u_j)\cap(K+u_j)\bigr)$
(in the appropriate nested form) is nonempty as soon as
$\sum_j\wid(P_j)/\ell_j<1$, which places the witness and forces
$\sum\wid(P_j)/\ell_j\ge1$ for coverings. The movement budget available in
direction $u$ is therefore governed by the homothety ratio $(\ell-a)/\ell$. This
ratio is tight: for a segment $K=[y,z]$ of length $\ell$ parallel to $u$,
$(K-u)\cap(K+u)$ is a segment of length exactly $\ell-a$. Hence the total budget in
direction $u$ is exactly the longest chord $\ell(u)$ and cannot be stretched to any
$N(u)>\ell(u)$; since $N_c(u)=(1-c)\ell(u)+c\wid_K(u)>\ell(u)$ whenever $c>0$ and
$\wid_K(u)>\ell(u)$, the claim follows.
\end{proof}
\begin{remark}
Proposition~\ref{thm:nogo} does \emph{not} show $\sum\wid(P_i)/N_c(u_i)\ge1$ is false
for $c>0$; it shows the sign-argument route cannot prove it. Beating $2/(1+\sqrt d)$ by
this normalizer therefore requires a witness that is not a Bang-system point, i.e.\
one exploiting the coupling of the covering. This is consistent with the
sharpness of~\cite{BY26} for the cube.
\end{remark}

\section{Open problems}
\begin{problem}
The three-plank case for three genuinely tilted directions on the triangle
(equivalently Conjecture~\ref{conj:bang} for the triangle). By
Theorem~\ref{thm:char} the remaining case is that of \emph{non-concurrent}
triples, where no single witness measure exists.
\end{problem}
\begin{problem}\label{prob:suff}
Theorem~\ref{thm:char} settles the existence of uniform-marginal witness
measures for three directions on the triangle. For $d\ge3$ the concurrence
condition of Proposition~\ref{prop:concur} remains necessary; is it sufficient
for $d+1$ directions on $\Delta^d$? The first test case is the cyclic medians
of $\Delta^d$, which concur at the centroid in every dimension
(\S\ref{sec:concur}).
\end{problem}
\begin{problem}
Whether $\sum\wid(P_i)/N_c(u_i)\ge1$ holds for some $c>0$ (a coupling certificate),
which would improve the general bound past $2/(1+\sqrt d)$.
\end{problem}

\appendix
\section{Edge tilings of cyclic triples}\label{app:tilings}

Throughout, $u_1=(\tau_1,0,1)$, $u_2=(0,1,\tau_2)$, $u_3=(1,\tau_3,0)$ is a
cyclic triple with arbitrary $\tau_1,\tau_2,\tau_3\in(0,1)$ (\emph{no}
concurrence is assumed), and $I_i=[l_i,h_i]\subset[0,1]$, $l_i<h_i$, are three
intervals. Say $(I_1,I_2,I_3)$ satisfies the \emph{edge-tiling conditions} if on
each edge of $\Delta$ the traces $T_i=\{x\in\text{edge}: u_i(x)\in I_i\}$ cover
the edge and have pairwise disjoint interiors. Indices are mod~$3$.

\begin{proposition}\label{prop:tilings}
For every $\tau\in(0,1)^3$ the edge-tiling conditions have exactly three
solutions: with
\[
\lambda_i=\frac{\tau_i\bigl(1-\tau_{i+1}(1-\tau_{i+2})\bigr)}{1+\tau_1\tau_2\tau_3},
\qquad
\rho_i=\frac{1-\tau_{i+2}(1-\tau_i)}{1+(1-\tau_1)(1-\tau_2)(1-\tau_3)},
\]
they are $\bigl([\lambda_i,\rho_i]\bigr)_i$, $\bigl([0,\lambda_i]\bigr)_i$, and
$\bigl([\rho_i,1]\bigr)_i$. One has $0<\lambda_i<\tau_i<\rho_i<1$ for all $i$ and
all $\tau$. For $\tau_i\equiv\tfrac12$ (the medians) the solutions are
$[\tfrac13,\tfrac23]^3$, $[0,\tfrac13]^3$, $[\tfrac23,1]^3$.
\end{proposition}

\subsection*{Setup}
Parametrize the edges by $t\in[0,1]$ in the orientations $A\!\to\!B$,
$B\!\to\!C$, $C\!\to\!A$. From the vertex values, each edge sees one direction
with values $[0,\tau_c]$ (the \emph{low} direction $c$, $u_c(t)=\tau_c(1-t)$),
one with values $[0,1]$ (the \emph{full} direction $f$, $u_f(t)=t$), and one
with values $[\tau_g,1]$ (the \emph{high} direction $g$,
$u_g(t)=1-(1-\tau_g)t$), with roles
\[
(c,f,g)=(1,2,3)\ \text{on } AB,\qquad (3,1,2)\ \text{on } BC,\qquad
(2,3,1)\ \text{on } CA .
\]
Inverting, the traces on an edge with roles $(c,f,g)$ are
\begin{equation}\label{eq:traces}
\begin{aligned}
T_c&=\Bigl[\max\bigl(0,1-\tfrac{h_c}{\tau_c}\bigr),\,1-\tfrac{l_c}{\tau_c}\Bigr]
\ \ (\text{nondegenerate iff }l_c<\tau_c),\qquad T_f=[l_f,h_f],\\
T_g&=\Bigl[\tfrac{1-h_g}{1-\tau_g},\,
\min\bigl(1,\tfrac{1-l_g}{1-\tau_g}\bigr)\Bigr]
\ \ (\text{nondegenerate iff }h_g>\tau_g).
\end{aligned}
\end{equation}
Classify each $I_i$ by its position relative to $\tau_i$: type $\mathrm L$ if
$h_i\le\tau_i$, type $\mathrm R$ if $l_i\ge\tau_i$, type $\mathrm M$ if
$l_i<\tau_i<h_i$; every $I_i$ is of exactly one type. Two elementary
observations, used throughout:

\emph{(O1)} If $I_c$ is of type $\mathrm M$ then $T_c=[0,\,1-l_c/\tau_c]$ (it
contains the left endpoint); if $I_c$ is of type $\mathrm R$ then $T_c$ is
degenerate (a point or empty). Dually, if $I_g$ is of type $\mathrm M$ then
$T_g=[(1-h_g)/(1-\tau_g),\,1]$; if $I_g$ is of type $\mathrm L$ then $T_g$ is
degenerate.

\emph{(O2)} The tiling condition is insensitive to degenerate traces: a
degenerate trace covers at most one point, and since the nondegenerate traces
are finitely many closed sets whose union contains $[0,1]$ minus finitely many
points, that union is all of $[0,1]$. So on each edge the \emph{nondegenerate}
traces tile $[0,1]$; in particular, finitely many closed intervals of positive
length with disjoint interiors covering $[0,1]$ abut consecutively from $0$
to~$1$.

\subsection*{Symmetries}
The rotation $A\to B\to C\to A$ maps the cyclic triple with parameters
$(\tau_1,\tau_2,\tau_3)$ to the cyclic triple with parameters
$(\tau_2,\tau_3,\tau_1)$ and acts on solutions by
$(I_1,I_2,I_3)\mapsto(I_2,I_3,I_1)$, shifting type vectors cyclically. The
reflection swapping $B$ and $C$, composed with the flip $s\mapsto1-s$ of each
normalized direction (which leaves every plank unchanged), maps the cyclic
triple $(\tau_1,\tau_2,\tau_3)$ to the cyclic triple
$(1-\tau_1,\,1-\tau_3,\,1-\tau_2)$ and acts by
$(I_1,I_2,I_3)\mapsto(1-I_1,\,1-I_3,\,1-I_2)$ (where $1-[l,h]=[1-h,1-l]$),
exchanging types $\mathrm L\leftrightarrow\mathrm R$. Since we prove
Proposition~\ref{prop:tilings} for \emph{all} $\tau$ at once, these symmetries
permute the class of problems, and the $27$ type vectors
$(X_1,X_2,X_3)\in\{\mathrm L,\mathrm M,\mathrm R\}^3$ reduce to seven orbit
representatives:
\[
\mathrm{MMM};\quad \mathrm{LLL}\ (\sim\mathrm{RRR});\quad \mathrm{MML};\quad
\mathrm{LLM};\quad \mathrm{LLR};\quad \mathrm{LMR};\quad \mathrm{LRM}
\]
(orbit sizes $1,2,6,6,6,3,3$). It suffices to analyze one representative per
orbit.

\subsection*{Case MMM}
By (O1), on every edge $T_c=[0,\,1-l_c/\tau_c]$ and
$T_g=[(1-h_g)/(1-\tau_g),\,1]$, both nondegenerate, and $T_f$ has positive
length. Disjointness of interiors forces $T_f$ inside the gap between them, and
covering forces the abutments
\[
l_f=1-\frac{l_c}{\tau_c},\qquad h_f=\frac{1-h_g}{1-\tau_g}.
\]
Reading this on the three edges gives two decoupled cyclic systems,
\[
l_2=1-\frac{l_1}{\tau_1},\quad l_1=1-\frac{l_3}{\tau_3},\quad
l_3=1-\frac{l_2}{\tau_2};
\qquad
h_1=\frac{1-h_2}{1-\tau_2},\quad h_2=\frac{1-h_3}{1-\tau_3},\quad
h_3=\frac{1-h_1}{1-\tau_1}.
\]
Each system composes three decreasing affine maps, so the composite self-map
(of $l_1$, resp.\ $h_1$) is affine with negative slope
($-1/(\tau_1\tau_2\tau_3)$, resp.\ $-1/\prod(1-\tau_i)$), hence has a unique
fixed point: the solutions are $l_i=\lambda_i$, $h_i=\rho_i$ (direct
substitution verifies the formulas; e.g.\
$1-\lambda_1/\tau_1
=\bigl(\tau_1\tau_2\tau_3+\tau_2(1-\tau_3)\bigr)/(1+\tau_1\tau_2\tau_3)
=\lambda_2$). The chain
$0<\lambda_i<\tau_i<\rho_i<1$ reduces to
$\tau_{i+1}(\tau_{i+2}-1)<\tau_1\tau_2\tau_3$ and $\tau_i(1-\tau_{i+1})<1$,
both automatic; so the solution is of type MMM and, by construction, tiles each
edge as $[0,l_f]\cup[l_f,h_f]\cup[h_f,1]$.

\subsection*{Case LLL}
All $T_g$ are degenerate, so each edge is tiled by
$T_c=[1-h_c/\tau_c,\,1-l_c/\tau_c]$ (the full image, as $I_c\subset[0,\tau_c]$)
and $T_f=I_f$ with $h_f\le\tau_f<1$. The tile containing $1$ must be $T_c$,
forcing $l_c=0$. If $h_c=\tau_c$ then $T_c=[0,1]$ and $T_f$ has empty interior,
impossible; so $T_f$ contains $0$ ($l_f=0$) and abuts $T_c$:
$h_f=1-h_c/\tau_c$. On the three edges: $l_i=0$ for all $i$, and the $h_i$
satisfy the \emph{same} cyclic system as the $l_i$ in case MMM, so
$h_i=\lambda_i$, giving $\bigl([0,\lambda_i]\bigr)_i$ (type $\mathrm L$ since
$\lambda_i<\tau_i$). Applying the reflection symmetry, the orbit representative
RRR yields the unique solution $h_i=1$, $l_i=\rho_i$, i.e.\
$\bigl([\rho_i,1]\bigr)_i$; equivalently, the $l_i$ satisfy the same cyclic
system as the $h_i$ in case MMM.

\subsection*{Case MML \textup(impossible\textup)}
On $AB$: $T_3$ is degenerate ($I_3$ of type $\mathrm L$), so $[0,1]$ is tiled
by $T_1=[0,\,1-l_1/\tau_1]$ (O1) and $T_2=I_2$. If $1\in T_1$ then $T_1=[0,1]$
and $T_2$ has empty interior, impossible; so $h_2=1$. Then on $BC$ the high
trace \eqref{eq:traces} is
$T_2=[\tfrac{1-h_2}{1-\tau_2},\min(1,\tfrac{1-l_2}{1-\tau_2})]=[0,1]$ (as
$h_2=1$ and $l_2<\tau_2$), leaving the full trace $T_1=I_1$ (positive length)
with empty interior --- contradiction.

\subsection*{Case LLM \textup(impossible\textup)}
On $BC$ ($(c,f,g)=(3,1,2)$): the high trace is degenerate ($I_2$ of type
$\mathrm L$), so $[0,1]$ is tiled by $T_3=[0,\,1-l_3/\tau_3]$ (O1, $I_3$ of
type $\mathrm M$) and $T_1=I_1$ with $h_1\le\tau_1<1$. The tile containing $1$
must be $T_3$, so $l_3=0$ and $T_3=[0,1]$, leaving $T_1$ with empty interior
--- contradiction.

\subsection*{Case LLR \textup(impossible\textup)}
On $BC$ both compressed traces are degenerate ($I_3$ of type $\mathrm R$ gives
a degenerate low trace, $I_2$ of type $\mathrm L$ a degenerate high trace), so
$T_1=I_1$ alone tiles $[0,1]$, forcing $h_1=1>\tau_1$ --- contradicting type
$\mathrm L$.

\subsection*{Case LMR \textup(impossible\textup)}
On $BC$: low trace degenerate ($I_3$ type $\mathrm R$), so $T_1=I_1$ and
$T_2=[\tfrac{1-h_2}{1-\tau_2},\,1]$ (O1) tile $[0,1]$; as above this forces
$l_1=0$ and $h_1=\tfrac{1-h_2}{1-\tau_2}$. On $CA$: high trace degenerate
($I_1$ type $\mathrm L$), so $T_2=[0,\,1-l_2/\tau_2]$ (O1) and $T_3=I_3$ tile,
forcing $h_3=1$ and $l_3=1-l_2/\tau_2$. Now on $AB$ the high trace is
$T_3=\bigl[0,\,\tfrac{1-l_3}{1-\tau_3}\bigr]
=\bigl[0,\,\tfrac{l_2}{\tau_2(1-\tau_3)}\bigr]$
(truncated at $1$ if larger, in which case $T_2$ already has empty interior, a
contradiction), of positive length, so $l_2>0$; since
$\tau_2(1-\tau_3)<1$, its right endpoint exceeds $l_2$, and the interiors of
$T_3$ and $T_2=[l_2,h_2]$ overlap --- contradiction.

\subsection*{Case LRM \textup(impossible\textup)}
On $CA$: both compressed traces are degenerate ($I_2$ of type $\mathrm R$,
$I_1$ of type $\mathrm L$), so $T_3=I_3$ alone tiles $[0,1]$, i.e.\
$I_3=[0,1]$. Then on $AB$ the high trace \eqref{eq:traces} is
$T_3=[\tfrac{1-h_3}{1-\tau_3},\min(1,\tfrac{1-l_3}{1-\tau_3})]=[0,1]$, leaving
$T_2=I_2$ (positive length) with empty interior --- contradiction.

\medskip
All orbits are exhausted, proving Proposition~\ref{prop:tilings}. \qed

\medskip
The three solutions have also been verified by exact rational searches at
$\tau=(\tfrac12,\tfrac12,\tfrac12)$ and $\tau=(\tfrac12,\tfrac23,\tfrac13)$,
and the median enumeration was verified by two further independent exact
computations; all scripts are archived as ancillary files.

\begin{thebibliography}{9}
\bibitem{Ambrus} G.\ Ambrus, \emph{Appendix: Plank problems}, manuscript, 2010.
\bibitem{Ball91} K.\ Ball, \emph{The plank problem for symmetric bodies},
Invent.\ Math.\ \textbf{104} (1991), 535--543.
\bibitem{Bang51} T.\ Bang, \emph{A solution of the ``plank problem''},
Proc.\ Amer.\ Math.\ Soc.\ \textbf{2} (1951), 990--993.
\bibitem{BY26} E.\ Bakaev, A.\ Yehudayoff, \emph{A note on the affine plank
conjecture}, arXiv:2602.20290, 2026.
\bibitem{Gardner88} R.\ J.\ Gardner, \emph{Relative width measures and the plank
problem}, Pacific J.\ Math.\ \textbf{135} (1988), no.\ 2, 299--312.
\bibitem{Verreault} W.\ Verreault, \emph{Plank theorems and their applications: a
survey}, arXiv:2203.05540v2, 2025.
\end{thebibliography}
\end{document}
